%% INDICON 2026 submission — DOUBLE-BLIND. No names, no affiliations, no
%% institutional identifiers. Hard limit 6 pages INCLUDING figures and references;
%% over six pages is an automatic reject.
%%
%% Compile on Overleaf with the IEEE Conference template (IEEEtran).
\documentclass[conference,a4paper,10pt]{IEEEtran}
\IEEEoverridecommandlockouts

\usepackage{cite}
\usepackage{amsmath,amssymb}
\usepackage{graphicx}
\usepackage{booktabs}
\usepackage{url}
\usepackage[hidelinks]{hyperref}

\begin{document}

\title{Reference-Set Substitution for Cross-Machine Pump Fault Diagnosis,\\
and What Honest Evaluation Costs}

%% DOUBLE-BLIND: leave exactly as-is for submission. Restore for camera-ready.
\author{\IEEEauthorblockN{Anonymous Authors}
\IEEEauthorblockA{Affiliation withheld for double-blind review}}

\maketitle

\begin{abstract}
Smallholder irrigation pumps fail silently, and the fastest failure is
destructive: under a minute of dry running ruins a mechanical seal. Per-pump
machine learning is impractical at this price point because commissioning cost
dominates hardware cost. We present a two-tier architecture in which
battery-powered nodes perform continuous statistical gating and a local dry-run
trip, while a shared gateway classifies escalated events with a prior-fitted
tabular foundation model, so that commissioning a new pump substitutes an
in-context reference set rather than retraining. On eleven in-service submersible
pumps under leave-one-machine-out, the approach reaches macro-F1
$0.738\pm0.015$ against a nested-tuned gradient-boosted baseline at $0.666$; at a
deployment-realistic budget of one false alarm per pump per month it recovers
$2.4\times$ as many faults. Accuracy saturates near 500 labelled reference
windows, giving a concrete commissioning specification, and a classification
costs 88\,ms on the target ARM gateway. We further report a protocol result:
under identical data and models, random-window splits nearly double reported
macro-F1 on in-service pumps and inflate it further on rig data, and the
inflation persists on both real datasets we examine. Finally we report that the gate's recall ceiling, not the
classifier, bounds the system: on the worst pump it is 0.52.
\end{abstract}

\begin{IEEEkeywords}
condition monitoring, fault diagnosis, centrifugal pumps, data leakage,
evaluation protocol, in-context learning, edge computing
\end{IEEEkeywords}

\section{Introduction}
%% ~500 words. Keep the two-timescale framing — it motivates the architecture.
An irrigation pump fails on two very different timescales, and conflating them is
the design error this work is organised around.

The failure that destroys equipment fastest is dry running. When suction is lost
the mechanical seal loses the fluid film that lubricates and cools it, and the
faces are damaged within tens of seconds. No architecture that ships a
measurement to a server and waits for a reply can prevent this. Dry running is
therefore not a classification problem in our design; it is a local trip.

The failures that dominate lifetime cost are slower: bearing wear, impeller
damage, misalignment, cavitation, a loosening foot. These develop over days to
months and are worth diagnosing rather than merely detecting, because the
maintenance action differs by fault. Here latency is irrelevant and
discrimination is everything.

The obstacle to machine learning at this price point is not inference cost but
\emph{commissioning} cost. A per-pump supervised model needs labelled examples of
each fault on each pump, which for a population of smallholder pumps costs more
than the pumps. We attack this with a prior-fitted tabular model that conditions
on a reference set supplied at inference time rather than fitting weights to it.
Commissioning becomes substituting a reference set, not retraining.

\textbf{Contributions.}
(C1) A two-tier architecture with an energy budget and gateway latency
\emph{measured on the deployment board}, and a gate escalation rate measured on
eleven in-service pumps.
(C2) Cross-machine evaluation of reference-set substitution under
leave-one-machine-out, against nested-tuned baselines at matched coverage.
(C3) An assessment of whether the expensive model earns its cost, at a
deployment-realistic alarm budget, with a commissioning specification.
(C4) A five-rung leakage ladder quantifying how much reported accuracy is
manufactured by the choice of split, and the observation that one widely used
pump dataset cannot support cross-machine evaluation at all.

We also report negative results that revise our own design: vibration is the
wrong primary sensor for dry running, ``training-free'' overstates what
in-context learning provides, and the edge accelerators on our deployment board
cannot take a model whose input shape varies by construction.

\section{Related Work}
%% ~450 words. The positioning against [6] is the load-bearing part — do not cut it.
Prior-data fitted networks invert the usual training relationship: a transformer
is pre-trained once on synthetic tabular tasks, then conditions on a labelled
reference set at inference time without gradient updates. Hollmann
\emph{et al.}~\cite{hollmann2025} show this outperforms tuned gradient-boosted
trees on datasets up to roughly ten thousand rows --- the regime a per-pump
commissioning set occupies. The property we need is not accuracy at small $n$ but
that adapting to a new machine is a data substitution rather than a training run.

The closest precedent for our classifier is Magad\'an \emph{et al.}~\cite{magadan2023},
who apply this model family to early fault classification in rotating machinery
under limited data. Three things separate our work: they evaluate the first
version of the model; their machines are motors and bearings rather than pumps,
so the taxonomy excludes cavitation, impeller damage and dry running; and their
evaluation is offline, with no deployment target, energy budget or gating tier.

Our protocol contribution sits in a small but growing literature arguing that
reported accuracy in vibration-based diagnosis is substantially an artefact of
how data are split. Varej\~ao \emph{et al.}~\cite{varejao2024} name the
phenomenon \emph{similarity bias} --- chunks of one chopped signal appearing on
both sides of the split --- and on the pump data we use report their best model
falling from an F-measure of 0.887 to 0.733 once their sampling strategy removes
it. Wheat \emph{et al.}~\cite{wheat2024} reach a compatible conclusion on bearing
data, distinguishing segmentation-level from bearing-level leakage.

Closest to our protocol contribution is Vieira \emph{et al.}~\cite{vieira2025},
who survey eighteen 2025 papers, find leakage persists in the majority, and
propose a methodology built on \emph{bearing-wise} partitioning across four
benchmarks. Their taxonomy and ours agree closely: their segment-wise and
repetition-wise splits are our levels 0--1, condition-wise is our level 3, and
bearing-wise --- the strictest they define --- is our level 2. The scenario they
motivate it with is a bearing replaced during maintenance on the same machine.

Our contribution is the rung above. Leave-one-machine-out holds out an entire
in-service pump: different installation, duty cycle and commissioning history,
not a swapped component on a shared rig. That is what a node deployed on a
customer's pump faces, and it is unavailable in the bearing benchmarks that
literature uses, because their components share a test rig. Vieira
\emph{et al.} also corroborate our central limitation independently, finding that
the number of unique training components decides robust performance --- the same
conclusion we reach about machines, on different data with a different model
family.

\section{System Design}
%% ~600 words + Fig. 1. The energy inversion is the most quotable result here.
\subsection{Two tiers, split by timescale}
The node is a battery-powered microcontroller with an accelerometer and a current
transformer. It runs a continuous low-cost statistic over the sensor stream,
trips the contactor locally on dry running, and escalates anomalous windows to
the gateway. It never classifies. The gateway is an ARM-class single-board
computer shared across a site; it holds the reference set and classifies
escalated feature windows. This split is a latency partition, not a capability
partition.

\subsection{The dry-run trip}
The trip requires three conditions \emph{simultaneously}: a CUSUM detector past
threshold, an absolute level below a floor, and persistence over a dwell period.
The conjunction is load-bearing. An earlier version disjoined the floor and the
CUSUM decision and tripped on essentially every startup transient --- a pump that
has not yet primed looks exactly like one that has lost suction for the few
seconds before it fills.

\textbf{Negative result.} The design began with a vibration-driven trip.
Vibration energy \emph{falls} when a pump loses suction, which is the opposite of
the signature a threshold detector wants, while motor current drops immediately
and unambiguously. We report this because the vibration-first design is the
intuitive one and it is wrong.

\subsection{The gate, and what actually costs energy}
The gate is commissioned on \emph{each pump's own healthy baseline}: a Mahalanobis
distance in the gate feature space plus an EWMA on level, with thresholds from
that pump's own healthy windows. This matches deployment --- a node is installed
on a presumed-healthy pump, observes it, and escalates what does not look like
what it saw. Commissioning requires $n > 10p$ healthy windows; estimated from
fewer, the covariance is degenerate and the distance rejects everything.

\textbf{The energy result inverts our starting assumption.} The design originally
budgeted radio transmission at roughly 55\,\% of the node's energy and optimised
the payload accordingly. That figure was computed for a \emph{fixed} transmission
schedule, which the same design had already rejected in favour of event-triggered
uplink. With a working gate, measured transmission is about 1\,\% of the node
budget and continuous sensing is essentially all of the rest. The optimisation
target is a cheaper always-on front end, not a smaller payload.

\subsection{Gateway inference}
The reference set is fixed between commissioning events, so the transformer's
key/value state can be cached once at boot rather than re-encoding the context on
every query. Measured on the deployment board --- a Rockchip RK3588, four
Cortex-A76 and four Cortex-A55 cores, single-threaded --- this is worth
$7.1\times$, and reducing the ensemble from eight members to one a further
$5.6\times$. Both ratios reproduce the values we measured on a development
workstation ($7.4\times$, $5.5\times$) to within 0.3, on hardware roughly
$21\times$ slower; the board is uniformly $21$--$22\times$ slower across all four
configurations with no configuration-dependent cliff, indicating a compute-bound
workload. In the deployed configuration a classification costs 88\,ms against two
to three escalations per pump per day, so latency is nowhere near binding.

\textbf{Negative result: no accelerator fits.} We did not attempt a port and do
not claim impossibility. Both the on-board NPU and an Edge TPU require a
fixed-shape, fully-INT8 graph from a restricted operator set. The reference set is
part of the model's input, so the input tensor has shape
$(n_\text{context}+n_\text{query}, n_\text{features})$ and varies with
commissioning. Padding would fix the shape but not the two obstacles that matter:
the operator sets do not cover a transformer attention stack, and INT8 quantising
a prior-fitted model without degrading the calibration our abstention depends on
is an open problem. The gateway is shared and mains-powered, so CPU inference
suffices.

\textbf{Negative result: ``training-free'' overstates it.} In-context learning
removes per-pump \emph{gradient} training. It does not remove the reference set,
the commissioning procedure, the feature pipeline, the per-pump gate thresholds
or the normalisation statistics.

\section{Datasets and Evaluation Protocol}
%% ~550 words + Table I. The Twente LOMO impossibility is a strong, citable claim.
\subsection{Two complementary real datasets}
No single public dataset supports this evaluation, so we use two that fail in
opposite directions (Table~\ref{tab:datasets}). ESPset~\cite{varejao2024} is
eleven in-service electrical submersible pumps with published order-normalised
velocity spectra, and is the only source of a genuine cross-machine axis. The
Twente/4TU dataset~\cite{kumar2023} is a laboratory rig: two motors at four
speeds, with both vibration and motor current.

We state a scope limitation plainly: ESPset pumps are offshore submersibles, not
smallholder irrigation pumps. They share the failure physics and nothing else,
including duty cycle, size and installation. The cross-machine result is evidence
about the method, not about irrigation pumps specifically.

\textbf{Twente cannot support leave-one-machine-out at all.} Its two motors share
only the healthy class: every fault present on one is absent from the other. Such
a fold would train on one label set and test on a disjoint one, producing a number
that looks like a cross-machine result and means nothing. Our splitter checks this
condition --- every fold must train on the classes it tests --- and marks folds
that fail it as uninterpretable rather than scoring them.

\begin{table}[t]
\caption{Two datasets, failing in opposite directions.}
\label{tab:datasets}
\centering
\small
\begin{tabular}{lcc}
\toprule
 & Twente/4TU & ESPset \\
\midrule
Machines & 2, disjoint faults & \textbf{11, shared} \\
Leave-one-machine-out & impossible & only source \\
Cross-operating-point & 50/75/100\,\% & order-normalised \\
Current channel & only source & --- \\
Waveforms & raw series & spectra only \\
\bottomrule
\end{tabular}
\end{table}

\subsection{The leakage ladder}
Our central protocol claim is that in machinery fault diagnosis the split, not the
model, determines most of the reported accuracy. We evaluate every model at five
levels of increasing strictness on identical features:
(0)~\emph{random-window} --- windows shuffled regardless of origin; \textbf{invalid},
since windows from one recording appear on both sides;
(1)~\emph{record-wise} --- no recording spans the split;
(2)~\emph{component-wise} --- no physical component instance spans the split;
(3)~\emph{cross-operating} --- train and test at different operating points;
(4)~\emph{leave-one-machine-out} --- the held-out machine is never seen.
Level 0 is reported not as a baseline but as an artefact, to quantify what common
practice buys.

\subsection{Normalisation is protocol, not preprocessing}
Per-machine normalisation is standard and quietly changes what is measured. We
evaluate two explicit strategies: \emph{per-machine}, where each machine is
standardised by its own statistics including the held-out one (transductive, and
legitimate when a node self-commissions on the target pump); and
\emph{train-pooled}, using training-machine statistics only (inductive, the
stricter reading). We report both and label every table, because on these data
the gap reaches 0.25 macro-F1 --- larger than the effect of the model choice.

\section{Results}
%% ~700 words + Tables II & III + Fig. 2. This is the section to protect if space runs short.
\subsection{Cross-machine classification}
Under leave-one-machine-out on eleven in-service pumps
(Table~\ref{tab:crossmachine}), reference-set substitution exceeds a nested-tuned
gradient-boosted baseline by $+0.072$ macro-F1 at matched coverage, roughly
$4.4\times$ the combined seed standard deviation. Baselines are tuned with
machine-grouped nested cross-validation, so no hyperparameter is selected using
the held-out machine. Tuning does not rescue them: logistic regression moves from
0.663 to 0.638 and the gradient-boosted model from 0.666 to 0.664.

We report abstaining and non-abstaining variants as separate models with separate
names. They are different systems, their coverage differs, and comparing an
abstaining model's accuracy against a full-coverage baseline is not a comparison.

\textbf{The qualification belongs with the result.} Per-machine confidence
intervals overlap. Seed noise is roughly 0.016 while the per-machine bootstrap
intervals span roughly 0.19. Machine count, not seed count, is the binding
statistical constraint; a twelfth pump would sharpen this far more than a sixth
seed.

\begin{table}[t]
\caption{Cross-machine (LOMO, 11 folds), train-pooled normalisation,
mean\,$\pm$\,s.d. over 5 seeds.}
\label{tab:crossmachine}
\centering
\small
\begin{tabular}{lcccc}
\toprule
Model & Macro-F1 & Acc. & Cov. & Per-machine CI \\
\midrule
Majority & $0.228\pm0.000$ & 0.837 & 1.00 & [0.290, 0.395] \\
Logistic & $0.663\pm0.000$ & 0.914 & 1.00 & [0.578, 0.787] \\
GBDT     & $0.666\pm0.006$ & 0.930 & 1.00 & [0.610, 0.779] \\
Ours (no abst.) & $\mathbf{0.738\pm0.015}$ & 0.911 & 1.00 & [0.583, 0.773] \\
Ours (abstaining) & $\mathbf{0.753\pm0.015}$ & 0.937 & 0.81 & [0.642, 0.810] \\
\bottomrule
\end{tabular}
\end{table}

\subsection{How much accuracy the split manufactures}
Table~\ref{tab:leakage} gives the headline protocol result: \textbf{on both real
datasets the invalid split roughly doubles the reported score.} On eleven
in-service pumps, holding out the machine costs 0.368 macro-F1 against a
random-window split of the same data with the same model; on the rig, holding out
the recording costs 0.501. A practitioner reporting the random-window number is
reporting close to twice what deployment will deliver. Our 1.9$\times$ on ESPset
is consistent in direction with the 1.21$\times$ that~\cite{varejao2024} report
on the same data from removing similarity bias alone, and larger, as the stricter
rung predicts.

\begin{table}[t]
\caption{Leakage inflation, identical data and models. Gradient-boosted baseline
shown; the effect holds for every model. Per-machine normalisation throughout.}
\label{tab:leakage}
\centering
\small
\begin{tabular}{lccc}
\toprule
Dataset & Random-window & Strictest valid & Inflation \\
\midrule
ESPset, 11 in-service pumps & 0.793 & 0.425 (LOMO) & $\mathbf{1.9\times}$ \\
Twente, 2-motor rig & 0.853 & 0.352 (record) & $\mathbf{2.4\times}$ \\
\bottomrule
\end{tabular}
\end{table}

\subsection{Does the expensive model earn its cost?}
Accuracy at a free choice of threshold is not the operational quantity; the
operational quantity is how much fault recall survives at an alarm rate an
operator tolerates. At one false alarm per pump per month --- 1080 decisions per
pump-month, a false alarm rate of $0.00093$ --- recall separates far more sharply
than macro-F1: majority 0.000, logistic 0.032, gradient-boosted 0.084, ours
\textbf{0.203}. The expensive model recovers $2.4\times$ as many faults as the
tuned baseline at the operating point that determines whether the system is
switched off. All pairwise differences are significant under an exact McNemar
test~\cite{dietterich1998}.

\textbf{Commissioning specification.} Sweeping reference-set size, accuracy
saturates near 500 labelled windows (macro-F1 0.739, 0.77\,s per query, against
0.712 at 250 and 0.672 at 100) and \emph{regresses} at 1000 (0.719, 1.22\,s). A
larger reference set is not simply better: past saturation it costs latency and
returns accuracy.

\subsection{The gate bounds the system}
End-to-end recall cannot exceed the gate's escalation recall, and that ceiling is
not uniform across pumps. Sweeping all 127 subsets of the seven deployable gate
features, the best achievable ceiling is \emph{flat} from three features upward
(0.865--0.873); what varies is the spread, with the best two-feature subset at
0.838 against a median of 0.467. Gate performance is therefore governed by which
features are chosen, not how many. Commissioning length binds only at seven
features, the first size at which a pump fails $n>10p$.

\textbf{The binding number is the worst machine, not the mean.} Across every
deployable configuration the best worst-machine ceiling is \textbf{0.52}, so the
two-tier system cannot exceed roughly 50\,\% recall on its worst pump however good
the gateway classifier is. We report the worst machine alongside the mean, and
the fault-count-pooled figure alongside both, because machines contributed
between 13 and 162 faulty windows each.

\section{Limitations}
\textbf{Machine count is binding.} Eleven machines is the largest cross-machine
axis in public pump data and is not enough for the per-machine intervals to
separate. Every other quantity we could increase is already past usefulness.
\textbf{The alarm budget bounds deployability}, and is the honest headline: 20\,\%
recall at one false alarm per pump per month makes this a triage aid, not a
replacement for inspection. \textbf{ESPset is offshore submersibles.} \textbf{Our
generic features underperform expert-designed ones}: no subset of the features our
extractor computes reaches the ceiling that the dataset's own published columns
achieve, which argues gate feature design needs per-fleet attention.
\textbf{No run-to-failure data}, so we make no remaining-useful-life claim;
detection stratified by the dataset's own severity grading is the honest
substitute. \textbf{No collected rig data}: the acquisition path is implemented
and exercisable end to end against a simulated backend, but no real acquisition
was performed and we claim no rig contribution.

\section{Conclusion}
Reference-set substitution transfers across machines: on eleven in-service pumps
never seen during commissioning, conditioning a prior-fitted tabular model on a
substituted reference set outperforms a nested-tuned gradient-boosted baseline,
and the advantage widens at a deployment-realistic alarm budget. The
commissioning cost that follows --- roughly 500 labelled windows, beyond which
more is not better --- is a concrete specification.

The architectural results are mostly corrections of our own assumptions. Dry
running belongs at the node and to the current channel. Transmission is not the
node's energy problem; continuous sensing is. And the gate, not the classifier,
bounds the system.

We expect the protocol result to outlast the model result. The model will be
superseded. The finding that random-window splits roughly double reported
macro-F1 on in-service data --- and that the inflation is \emph{smallest} on the
simulated data where practitioners are most likely to look for it --- applies to
any model anyone puts on this problem next.

%% DOUBLE-BLIND: no acknowledgements, no funding, no artifact URL that identifies you.
\section*{Reproducibility}
Neither dataset is redistributed. ESPset is available from Mendeley Data under
CC BY 4.0 (version 3, not the version 1 that some prior work cites) and the
Twente/4TU dataset from 4TU.ResearchData under CC BY 4.0. Source code, the
evaluation harness and the result files behind every table are available to
reviewers as anonymised supplementary material and will be released publicly on
acceptance. This work is Built with PriorLabs-TabPFN.

\begin{thebibliography}{99}
\bibitem{hollmann2025} N.~Hollmann \emph{et al.}, ``Accurate predictions on small
data with a tabular foundation model,'' \emph{Nature}, vol.~637, pp.~319--326, 2025.
\bibitem{magadan2023} L.~Magad\'an, J.~Rold\'an-G\'omez, J.~C.~Granda, and
F.~J.~Su\'arez, ``Early fault classification in rotating machinery with limited
data using TabPFN,'' \emph{IEEE Sensors J.}, vol.~23, no.~24, pp.~30960--30970, 2023.
%% TODO verify end page
\bibitem{varejao2024} I.~M.~S.~Varej\~ao \emph{et al.}, ``An open source
experimental framework and public dataset for vibration-based fault diagnosis of
electrical submersible pumps used on offshore oil exploration,''
\emph{Knowledge-Based Systems}, vol.~289, art.~111452, 2024.
\bibitem{wheat2024} L.~Wheat, M.~von Mohrenschildt, S.~Habibi, and D.~Al-Ani,
``Impact of data leakage in vibration signals used for bearing fault diagnosis,''
\emph{IEEE Access}, vol.~12, pp.~169879--169895, 2024.
\bibitem{vieira2025} J.~P.~Vieira, V.~A.~Bauler, R.~K.~Rosa, and D.~Silva,
``Towards a more realistic evaluation of machine learning models for bearing
fault diagnosis,'' \emph{arXiv:2509.22267}, 2025.
%% TODO check for a journal version before submitting
\bibitem{kumar2023} D.~Kumar \emph{et al.}, ``Motor current and vibration
monitoring dataset for various faults in an E-motor-driven centrifugal pump,''
\emph{Data in Brief}, vol.~51, art.~109779, 2023.
\bibitem{dietterich1998} T.~G.~Dietterich, ``Approximate statistical tests for
comparing supervised classification learning algorithms,'' \emph{Neural
Computation}, vol.~10, no.~7, pp.~1895--1923, 1998.
\bibitem{demsar2006} J.~Dem\v{s}ar, ``Statistical comparisons of classifiers over
multiple data sets,'' \emph{JMLR}, vol.~7, pp.~1--30, 2006.
%% Add 8-12 more: pump fault physics, MCSA, CUSUM, LoRa energy, selective
%% prediction. INDICON reviewers expect a fuller reference list than eight.
\end{thebibliography}

\end{document}
